\documentclass[11pt,a4paper]{article}
\usepackage[T1]{fontenc}
\usepackage[utf8]{inputenc}
\usepackage[french]{babel}
\usepackage{geometry}
\usepackage{hyperref}
\usepackage{enumitem}
\usepackage{booktabs}
\geometry{margin=2.5cm}

\title{Rapport Technique -- Projet PD Classification (BFPME)}
\author{}
\date{}

\begin{document}
\maketitle

\renewcommand{\contentsname}{Sommaire}
\setcounter{tocdepth}{2}
\tableofcontents
\newpage

\section{Contexte et objectif du projet}

Ce projet vise a construire une base de travail pour la prediction du risque de defaut (PD) dans un cadre de modernisation progressive du scoring BFPME.

L'approche retenue est prudente:
\begin{itemize}[leftmargin=1.2cm]
  \item conserver la logique metier existante BFPME comme reference,
  \item creer un dataset centralise enrichi,
  \item tester des modeles de classification ML (RandomForest, XGBoost),
  \item comparer plusieurs strategies d'integration du signal BFPME.
\end{itemize}

Le but n'est pas de remplacer immediatement le scorecard historique, mais de demontrer une valeur ajoutee methodologique dans un cadre champion/challenger.

\section{Creation de \texttt{data/dataset\_final.csv}}

\subsection{Source de depart}

Le fichier \texttt{data/dataset\_final.csv} a ete construit a partir de:
\begin{itemize}[leftmargin=1.2cm]
  \item \texttt{data/bfpme\_synthetic\_dataset.csv} (base initiale contenant sous-facteurs BFPME, score et classes de risque),
  \item script de generation: \texttt{scripts/generate\_dataset\_final.py}.
\end{itemize}

La base finale contient:
\begin{itemize}[leftmargin=1.2cm]
  \item les variables d'origine BFPME (sous-facteurs + scores),
  \item des variables synthetiques additionnelles (business, credit, finance, garanties, comportement bancaire, macro),
  \item des ratios derives,
  \item une cible de classification binaire: \texttt{default\_flag}.
\end{itemize}

\subsection{Structure du dataset final}

Le dataset final est centralise ("wide table"):
\begin{itemize}[leftmargin=1.2cm]
  \item 1 ligne = 1 dossier client,
  \item 1 colonne = 1 variable explicative ou cible.
\end{itemize}

Il contient notamment:
\begin{itemize}[leftmargin=1.2cm]
  \item identifiants et metadonnees (\texttt{dossier\_id}, \texttt{code\_modele}),
  \item variables BFPME (\texttt{score\_bfpme\_*}, \texttt{risk\_class}, \texttt{SF\_*}),
  \item variables metier enrichies (juridique, taille, gouvernance, localisation, credit),
  \item variables financieres (structure bilan, rentabilite, liquidite, cashflow),
  \item variables comportementales et relation bancaire (DPD, impayes, discipline),
  \item variables macroeconomiques (\texttt{scenario\_macro}, inflation, PIB, etc.),
  \item cible: \texttt{default\_flag} (0/1) + variable intermediaire \texttt{pd\_target}.
\end{itemize}

\section{Logique de generation metier et formule latente de la cible}

\subsection{Principe general}

Le dataset a ete genere avec une logique de coherence metier, et non par tirage totalement aleatoire.

Une variable cachee est introduite:
\begin{itemize}[leftmargin=1.2cm]
  \item \texttt{latent\_risk} (niveau de risque latent),
  \item \texttt{latent\_quality = 1 - latent\_risk}.
\end{itemize}

Cette variable gouverne ensuite de nombreuses distributions:
\begin{itemize}[leftmargin=1.2cm]
  \item profils solides $\rightarrow$ meilleure liquidite, meilleure gouvernance, meilleur comportement de paiement,
  \item profils fragiles $\rightarrow$ endettement plus eleve, DPD plus fort, stress bancaire plus marque.
\end{itemize}

\subsection{Construction du risque latent}

Le risque latent est calcule a partir de:
\begin{itemize}[leftmargin=1.2cm]
  \item score BFPME global normalise,
  \item index de risque derive de la classe BFPME (\texttt{bfpme\_risk\_index}),
  \item type de modele (\texttt{code\_modele}),
  \item bruit gaussien modere.
\end{itemize}

Forme simplifiee:
\begin{quote}
\texttt{latent\_risk ~= 0.55 * base\_risk + 0.35 * bfpme\_risk\_index + 0.10 * I(code\_modele=1) + bruit}
\end{quote}

avec:
\begin{itemize}[leftmargin=1.2cm]
  \item \texttt{base\_risk = 1 - score\_bfpme\_global / 200}.
\end{itemize}

\subsection{Generation de \texttt{pd\_target}}

La probabilite de defaut intermediaire est obtenue via un score latent de defaut, combine avec une fonction logistique (sigmoid).

Le score latent de defaut agrege:

\textbf{Effets qui augmentent le risque}:
\begin{itemize}[leftmargin=1.2cm]
  \item faible score BFPME,
  \item DPD moyen eleve,
  \item leverage eleve (\texttt{debt\_to\_equity\_ratio}),
  \item faible capacite de service de la dette (\texttt{dscr} faible),
  \item conditions macro defavorables,
  \item comportements de stress (augmentation dettes, baisse depots, etc.).
\end{itemize}

\textbf{Effets qui reduisent le risque}:
\begin{itemize}[leftmargin=1.2cm]
  \item meilleure liquidite (\texttt{ratio\_liquidite\_generale}),
  \item meilleure couverture de garanties (\texttt{garantie\_coverage\_ratio}),
  \item meilleur taux de remboursement ponctuel,
  \item meilleure confiance institutionnelle.
\end{itemize}

Forme simplifiee:
\begin{quote}
\texttt{pd\_target = sigmoid(score\_latent\_default)}
\end{quote}

\subsection{Construction de la cible binaire \texttt{default\_flag}}

La cible finale de classification est un tirage Bernoulli:
\begin{quote}
\texttt{default\_flag ~ Bernoulli(pd\_target)}
\end{quote}

Donc:
\begin{itemize}[leftmargin=1.2cm]
  \item \texttt{default\_flag = 1}: defaut,
  \item \texttt{default\_flag = 0}: non-defaut.
\end{itemize}

Cette approche permet de transformer une logique de probabilite en probleme supervise binaire.

\section{Data preprocessing \& feature engineering}

Notebook de reference:
\begin{itemize}[leftmargin=1.2cm]
  \item \texttt{notebook/data\_preprocessing\_feature\_engineering.ipynb}
\end{itemize}

\subsection{Etapes de preprocessing appliquees}

\begin{enumerate}[leftmargin=1.2cm, label=\arabic*., widest=9]
  \item Chargement et controles initiaux:
  \begin{itemize}
    \item dimensions du dataset,
    \item types,
    \item taux de valeurs manquantes,
    \item distribution de la cible.
  \end{itemize}

  \item Controle de fuite d'information:
  \begin{itemize}
    \item suppression de \texttt{pd\_target} (utilisee pour generer la cible),
    \item suppression de l'identifiant \texttt{dossier\_id}.
  \end{itemize}

  \item Analyse exploratoire ciblee:
  \begin{itemize}
    \item visualisation de \texttt{default\_flag},
    \item distributions/boxplots de variables cle (score BFPME, DPD, liquidite, garanties, leverage, DSCR).
  \end{itemize}

  \item Detection des outliers:
  \begin{itemize}
    \item methode IQR sur variables cle.
  \end{itemize}

  \item Winsorization ciblee:
  \begin{itemize}
    \item clipping quantiles (1\% / 99\%) sur variables financieres extremes.
  \end{itemize}

  \item Encodage categoriel controle:
  \begin{itemize}
    \item \texttt{OrdinalEncoder} pour categories ordonnees,
    \item \texttt{OneHotEncoder} limite aux variables nominales necessaires.
  \end{itemize}

  \item Standardisation:
  \begin{itemize}
    \item \texttt{StandardScaler} sur variables numeriques.
  \end{itemize}

  \item Split train/test:
  \begin{itemize}
    \item split stratifie (80/20) sur \texttt{default\_flag}.
  \end{itemize}

  \item Export:
  \begin{itemize}
    \item jeux preprocesses (\texttt{X\_train\_preprocessed}, \texttt{X\_test\_preprocessed}, \texttt{y\_train}, \texttt{y\_test}).
  \end{itemize}
\end{enumerate}

\subsection{Feature engineering metier ajoute}

Variables composites principales:
\begin{itemize}[leftmargin=1.2cm]
  \item \texttt{debt\_total = dettes\_court\_terme + dettes\_long\_terme},
  \item \texttt{loan\_to\_revenue = montant\_pret / chiffre\_affaires\_annuel},
  \item \texttt{cash\_to\_debt = cash\_equivalents / debt\_total},
  \item \texttt{impayes\_to\_loan = montant\_impayes / montant\_pret},
  \item \texttt{stress\_behavior\_index} (somme de signaux de stress bancaire),
  \item \texttt{governance\_risk\_index} (experience + turnover + structure dirigeante).
\end{itemize}

Ces variables visent a mieux capter:
\begin{itemize}[leftmargin=1.2cm]
  \item pression de dette,
  \item fragilite de tresorerie,
  \item stress comportemental,
  \item risque de gouvernance.
\end{itemize}

\subsection{Realisme de la cible synthetique et flux \texttt{dataset\_final\_realism\_v2.csv}}

En complement des etapes ci-dessus, le notebook \texttt{notebook/data\_preprocessing\_feature\_engineering.ipynb} inclut une section intitulee \emph{Making the synthetic target more realistic} (en debut de cette partie du notebook, les objectifs sont resumes en quatre points: stochasticite accrue, reduction du poids des predicteurs dominants, contradictions realistes, split temporel).

Concretement, ce bloc du notebook:
\begin{enumerate}[leftmargin=1.2cm, label=\arabic*., widest=4]
  \item augmente le bruit dans la generation de la cible pour reduire le caractere quasi-deterministe de la premiere version;
  \item diminue la dominance relative de variables telles que \texttt{score\_bfpme\_global} et le DPD en repartissant le signal sur d'autres dimensions (macro, levier, DSCR, stress, protections);
  \item introduit des cas contradictoires (bons dossiers en defaut, dossiers risques qui tiennent);
  \item remplace le split aleatoire par un decoupage temporel (historique d'apprentissage, periode de test plus recente).
\end{enumerate}

Le resultat est exporte sous \texttt{data/dataset\_final\_realism\_v2.csv}, avec une cible de modelisation \texttt{default\_flag\_v2} et une probabilite intermediaire \texttt{pd\_target\_v2}. Les notebooks de classification (options A, B et C) utilisent ce meme jeu et la meme logique d'entree (chargement de \texttt{dataset\_final\_realism\_v2.csv}, cible \texttt{default\_flag\_v2}, et reprise des exports temporels \texttt{X\_train\_temporal\_v2} / \texttt{X\_test\_temporal\_v2} lorsque applicable) afin d'observer des metriques plus credibles que sur la premiere generation synthetique.

\section{Strategie de modelisation: 3 options comparees}

Trois notebooks de classification ont ete crees:
\begin{itemize}[leftmargin=1.2cm]
  \item \texttt{notebook/classification\_option\_A\_all\_features.ipynb}
  \item \texttt{notebook/classification\_option\_B\_without\_bfpme.ipynb}
  \item \texttt{notebook/classification\_option\_C\_hybrid\_min\_bfpme.ipynb}
\end{itemize}

Chaque notebook entraine:
\begin{itemize}[leftmargin=1.2cm]
  \item \texttt{RandomForestClassifier}
  \item \texttt{XGBoost} (ou fallback GradientBoosting si xgboost indisponible)
\end{itemize}

\subsection{Option A -- All Features}
\textbf{Objectif}
\begin{itemize}[leftmargin=1.2cm]
  \item mesurer la performance maximale avec toutes les variables utiles.
\end{itemize}
\textbf{Contenu}
\begin{itemize}[leftmargin=1.2cm]
  \item variables enrichies + signaux BFPME complets (scores et classe).
\end{itemize}
\textbf{Lecture metier}
\begin{itemize}[leftmargin=1.2cm]
  \item benchmark "upper bound" de performance.
\end{itemize}

\subsection{Option B -- Without BFPME Score Family}
\textbf{Objectif}
\begin{itemize}[leftmargin=1.2cm]
  \item evaluer la capacite predictive d'un modele plus "data-driven", moins dependant du score historique.
\end{itemize}
\textbf{Contenu}
\begin{itemize}[leftmargin=1.2cm]
  \item retrait des variables de score BFPME et classes associees.
\end{itemize}
\textbf{Lecture metier}
\begin{itemize}[leftmargin=1.2cm]
  \item mesure de la valeur predictive des nouvelles variables hors scorecard legacy.
\end{itemize}

\subsection{Option C -- Hybrid Minimal BFPME}
\textbf{Objectif}
\begin{itemize}[leftmargin=1.2cm]
  \item tester une approche hybride parcimonieuse.
\end{itemize}
\textbf{Contenu}
\begin{itemize}[leftmargin=1.2cm]
  \item conservation d'un seul signal BFPME (\texttt{score\_bfpme\_global}), retrait du reste de la famille score detaillee.
\end{itemize}
\textbf{Lecture metier}
\begin{itemize}[leftmargin=1.2cm]
  \item compromis entre explicabilite metier historique et signal data-driven.
\end{itemize}

\section{Resultats observes et interpretation}

Deux regimes d'evaluation ont ete utilises.

\subsection{Premiere generation (\texttt{dataset\_final.csv}, cible \texttt{default\_flag})}

Les metriques etaient tres elevees (accuracy/F1/AUC souvent autour de 0.96--0.99) sur les trois options, avec des ecarts modestes entre RandomForest et XGBoost.

\subsection{Generation realiste (\texttt{dataset\_final\_realism\_v2.csv}, cible \texttt{default\_flag\_v2})}

Apres application des quatre leviers de realisme (section~4.3) et evaluation sur \textbf{split temporel} (exports \texttt{X\_train\_temporal\_v2}, \texttt{X\_test\_temporal\_v2}, \texttt{y\_*\_temporal\_v2}), les metriques se situent dans une fourchette \textbf{plus credible pour un cas d'usage risque} qu'une PD synthetique quasi parfaite.

\paragraph{1) Pourquoi les scores etaient trop eleves avant.}
Dans la premiere version synthetique, la cible etait fortement liee a quelques predicteurs dominants (\texttt{score\_bfpme\_global}, variables de comportement type DPD), avec peu de bruit et un split aleatoire train/test tres proche en distribution. Les modeles pouvaient alors retrouver presque directement la logique de generation, d'ou des AUC proches de~1 (sur-ajustement au generateur plutot qu'une generalisation realiste).

\paragraph{2) Ce qui a change en v2 et impact attendu.}
\begin{itemize}[leftmargin=1.2cm]
  \item \textbf{Stochasticite accrue}: ajout d'une composante aleatoire plus forte dans la cible, ce qui augmente l'incertitude irreductible.
  \item \textbf{Moindre dominance des predicteurs top}: reduction du poids relatif de \texttt{score\_bfpme\_global} et du DPD, avec redistribution du signal vers macro, levier, DSCR, stress et protections.
  \item \textbf{Contradictions realistes}: injection de cas contre-intuitifs (bons dossiers qui defautent, dossiers risques qui tiennent).
  \item \textbf{Split temporel}: train sur dossiers plus anciens, test sur dossiers plus recents (derive de distribution plus proche du deploiement).
\end{itemize}

Ces quatre ajustements cassent les raccourcis deterministes et rendent l'evaluation plus exigeante.

\paragraph{3) Metriques observees (RandomForest et XGBoost, jeu v2 + split temporel).}
Les tableaux suivants reprennent les sorties des notebooks d'options A, B et~C sur le meme protocole (\texttt{dataset\_final\_realism\_v2.csv}, exports temporels). Les resultats \textbf{different entre A et B}: l'option~A conserve toute la famille score BFPME disponible, tandis que l'option~B l'exclut.

\begin{table}[h]
  \centering
  \caption{Option A -- toutes les variables utiles (hors fuite), \texttt{dataset\_final\_realism\_v2.csv}, split temporel (sortie notebook).}
  \begin{tabular}{lrrrrr}
    \toprule
    Modele & Accuracy & Precision & Recall & F1 & AUC \\
    \midrule
    RandomForest & 0.765 & 0.829060 & 0.782258 & 0.804979 & 0.819503 \\
    XGBoost      & 0.745 & 0.782946 & 0.814516 & 0.798419 & 0.802101 \\
    \bottomrule
  \end{tabular}
\end{table}

\begin{table}[h]
  \centering
  \caption{Option B -- sans variables de score BFPME (famille score retiree), meme jeu et split temporel que l'option~A (sortie notebook).}
  \begin{tabular}{lrrrrr}
    \toprule
    Modele & Accuracy & Precision & Recall & F1 & AUC \\
    \midrule
    RandomForest & 0.765 & 0.829060 & 0.782258 & 0.804979 & 0.820458 \\
    XGBoost      & 0.760 & 0.801587 & 0.814516 & 0.808000 & 0.793081 \\
    \bottomrule
  \end{tabular}
\end{table}

\begin{table}[h]
  \centering
  \caption{Option C -- hybrid minimal (\texttt{score\_bfpme\_global} conserve, reste de la famille score retire), meme protocole.}
  \begin{tabular}{lrrrrr}
    \toprule
    Modele & Accuracy & Precision & Recall & F1 & AUC \\
    \midrule
    RandomForest & 0.770 & 0.830508 & 0.790323 & 0.809917 & 0.820140 \\
    XGBoost      & 0.785 & 0.813953 & 0.846774 & 0.830040 & 0.811969 \\
    \bottomrule
  \end{tabular}
\end{table}

Les AUC autour de \textbf{0,79--0,82} et des scores F1 proches de \textbf{0,80--0,83} sont coherents avec une discrimination utile mais non ``parfaite''; ils sont preferables, pour la communication interne et la preparation champion/challenger, a des AUC $>$~0,95 issues d'une cible quasi rejouable.

\paragraph{3.b) Comparaison de robustesse: 1000 vs 10000 clients (meme logique v2).}
En augmentant la taille de la population synthetique (de 1000 a 10000 clients) tout en conservant la meme logique de generation et les memes colonnes, les metriques baissent de facon homogene et convergent vers un plateau plus strict.

\begin{table}[h]
  \centering
  \caption{Resultats sur 10000 clients (dataset v3, split temporel, memes options A/B/C).}
  \begin{tabular}{llrrrrr}
    \toprule
    Option & Modele & Accuracy & Precision & Recall & F1 & AUC \\
    \midrule
    A & RandomForest & 0.733600 & 0.783307 & 0.788368 & 0.785829 & 0.754300 \\
    A & XGBoost      & 0.729594 & 0.780096 & 0.785137 & 0.782609 & 0.751140 \\
    B & RandomForest & 0.735103 & 0.783373 & 0.791599 & 0.787465 & 0.753904 \\
    B & XGBoost      & 0.729094 & 0.777689 & 0.788368 & 0.782992 & 0.752760 \\
    C & RandomForest & 0.734602 & 0.783200 & 0.790792 & 0.786977 & 0.753921 \\
    C & XGBoost      & 0.728593 & 0.781553 & 0.780291 & 0.780922 & 0.752605 \\
    \bottomrule
  \end{tabular}
\end{table}

Lecture comparative avec les resultats ``1000 clients'' ci-dessus:
\begin{itemize}[leftmargin=1.2cm]
  \item l'AUC passe globalement d'une zone \textbf{0,79--0,82} vers \textbf{0,75--0,76} (baisse d'environ 4 a 7 points selon option/modele);
  \item accuracy et F1 diminuent egalement (ordre de grandeur proche de 5 a 6 points), avec un comportement relativement stable entre A/B/C;
  \item les ecarts inter-options deviennent faibles sur 10000 clients, ce qui suggere que la difficulte provient davantage de la cible realiste et du split temporel que d'un effet artificiel de petite taille d'echantillon;
  \item ces niveaux restent coherents avec une discrimination utile dans un cadre risque plus exigeant, et renforcent le caractere defendable de l'approche avant passage aux donnees reelles.
\end{itemize}

\paragraph{4) Lecture metier pour presentation BFPME/client.}
\begin{itemize}[leftmargin=1.2cm]
  \item v1 a surtout valide la coherence technique du pipeline;
  \item v2 valide une robustesse plus credible sous hypotheses realistes (bruit, contradictions, derive temporelle);
  \item la performance reste solide apres stress de realisme, ce qui rend la methode plus defendable avant migration sur donnees reelles.
\end{itemize}

\paragraph{5) Etape pratique recommandee.}
Utiliser par defaut les exports temporels v2 comme base commune de comparaison A/B/C:
\begin{itemize}[leftmargin=1.2cm]
  \item \texttt{data/X\_train\_temporal\_v2.csv},
  \item \texttt{data/X\_test\_temporal\_v2.csv},
  \item \texttt{data/y\_train\_temporal\_v2.csv},
  \item \texttt{data/y\_test\_temporal\_v2.csv}.
\end{itemize}

\section{Livrables techniques existants}

\begin{itemize}[leftmargin=1.2cm]
  \item Dataset final (generation initiale):
  \begin{itemize}
    \item \texttt{data/dataset\_final.csv}
  \end{itemize}

  \item Dataset synthetique realiste (cible \texttt{default\_flag\_v2}):
  \begin{itemize}
    \item \texttt{data/dataset\_final\_realism\_v2.csv}
  \end{itemize}

  \item Exports temporels et preprocesses (v2, selon le notebook):
  \begin{itemize}
    \item \texttt{data/X\_train\_temporal\_v2.csv}, \texttt{data/X\_test\_temporal\_v2.csv}
    \item \texttt{data/y\_train\_temporal\_v2.csv}, \texttt{data/y\_test\_temporal\_v2.csv}
    \item \texttt{data/X\_train\_preprocessed\_temporal\_v2.csv}, \texttt{data/X\_test\_preprocessed\_temporal\_v2.csv} (si generes par le flux ordonne du notebook)
  \end{itemize}

  \item Documentation variables:
  \begin{itemize}
    \item \texttt{dataset\_final\_variable\_dictionary.txt}
  \end{itemize}

  \item Generation du dataset:
  \begin{itemize}
    \item \texttt{scripts/generate\_dataset\_final.py}
  \end{itemize}

  \item Preprocessing \& feature engineering:
  \begin{itemize}
    \item \texttt{notebook/data\_preprocessing\_feature\_engineering.ipynb}
  \end{itemize}

  \item Classification (3 options):
  \begin{itemize}
    \item \texttt{notebook/classification\_option\_A\_all\_features.ipynb}
    \item \texttt{notebook/classification\_option\_B\_without\_bfpme.ipynb}
    \item \texttt{notebook/classification\_option\_C\_hybrid\_min\_bfpme.ipynb}
  \end{itemize}
\end{itemize}

\section{Conclusion}

Le projet fournit un cadre complet et presentable de modelisation PD en mode "proof of value":
\begin{itemize}[leftmargin=1.2cm]
  \item conception de donnees enrichies,
  \item cible de classification binaire,
  \item pipeline reproductible,
  \item comparaison strategique A/B/C.
\end{itemize}

La premiere generation synthetique produisait des metriques volontairement tres elevees. La seconde (\texttt{dataset\_final\_realism\_v2.csv}, avec bruit, moindre dominance des predicteurs dominants, contradictions et split temporel) ramene les AUC vers environ \textbf{0,79--0,82}: ce niveau est plus \textbf{defendable pour le cas d'usage} car il limite l'illusion d'une PD ``rejouable'' a partir des seules entrees. L'ensemble reste synthetique: la validation definitive repose sur des donnees reelles avec outcomes observes.

Cette demarche reste tres utile pour:
\begin{itemize}[leftmargin=1.2cm]
  \item industrialiser la methodologie,
  \item aligner les equipes risque/data,
  \item preparer la transition vers un challenger modele sur donnees reelles BFPME.
\end{itemize}

\end{document}

